% Figure for Oscillations, Waves and Optics §A4.2 (a corner reflector: two perpendicular mirrors; each reflection reverses one component of the ray direction, so the outgoing ray is antiparallel to the incoming one).
\begin{tikzpicture}[font=\small, >={Stealth[length=2.4mm]},
    mid/.style={postaction={decorate}, decoration={markings, mark=at position 0.5 with {\arrow{Stealth[length=2.6mm]}}}}]
  % mirrors with hatching on the back side
  \fill[pattern=north east lines, pattern color=black!40] (-0.25,-0.25) rectangle (4.0,0);
  \fill[pattern=north east lines, pattern color=black!40] (-0.25,0) rectangle (0,3.8);
  \draw[black!75, very thick] (4.0,0) -- (0,0) -- (0,3.8);
  \node[black!60, font=\scriptsize, anchor=north] at (2.6,-0.3) {mirror 1};
  \node[black!60, font=\scriptsize, rotate=90, anchor=south] at (-0.3,2.7) {mirror 2};
  \def\c{0.8192}\def\s{0.5736}
  \coordinate (S) at (3.6,1.4); \coordinate (P1) at (1.6,0); \coordinate (P2) at (0,1.1209);
  \draw[figB, very thick, mid] (S) -- (P1);
  \draw[figB, very thick, mid] (P1) -- (P2);
  \draw[figR, very thick, mid] (P2) -- ($(P2)+4.3*(\c,\s)$);
  \node[figB, font=\scriptsize, anchor=west] at (S) {in: $(-a,\,-b)$};
  \node[figR, font=\scriptsize, anchor=west] at ($(P2)+4.3*(\c,\s)$) {out: $(+a,\,+b)$};
  % normals at the reflection points
  \draw[black!45, dashed] (1.6,0) -- (1.6,1.2);
  \draw[black!45, dashed] (0,1.1209) -- (1.3,1.1209);
\end{tikzpicture}
